\chapter{تنسيق وإشارات التزامن عبر النوم واليقظة}
\label{CH:SLEEP}

تحتاج العمليات داخل النواة غالبًا إلى الانتظار لحدوث حدث معين (مثل قراءة بيانات من القرص الصلب، انتظار إدخال من الشاشة، أو انتظار انتهاء عملية ابن عبر \texttt{wait()}).

إذا انتظرت العملية باستخدام حلقة تكرارية نشطة (Spin-waiting)، فإنها تستنزف موارد المعالج دون جدوى.
يوفر \texttt{xv6} آلية التنسيق \textbf{النوم واليقظة} (Sleep and Wakeup - أو أسلوب الحراس والشرطية)، التي تسمح للعملية بالتعليق والدخول في حالة النوم حتى يقوم جزء آخر من النواة بإيقاظها عند تحقق الشرط.

يقدم هذا الفصل آليات النوم واليقظة، مشكلة الاستيقاظ المفقود (Lost Wakeup)، وكيفية تنفيذ الأقفال النوامة، واستدعاءات النظام \texttt{wait()} و \texttt{exit()} و \texttt{kill()}.

\section{نظرة عامة}
\label{sec:sleep_overview}

تعتمد الآلية البسيطة للتنسيق على وجود عنوان انتسساب أو قناة انتظار تُدعى \textbf{قناة النوم} (chan / sleep channel):
\begin{itemize}
  \item \texttt{sleep(chan, lock)}: تضع العملية الحالية في حالة النوم على العنوان \texttt{chan} وتتخلى عن المعالج.
  \item \texttt{wakeup(chan)}: توقظ كافة العمليات المنتظرة أو النائمة على العنوان \texttt{chan} وتغير حالتها إلى \texttt{RUNNABLE}.
\end{itemize}

لتجنب مشكلة \textbf{الاستيقاظ المفقود} (Lost Wakeup):
تحدث هذه المشكلة إذا قُطعت العملية أثناء فحص الشرط وقبل النوم مباشرةً، فتقوم دالة المقاطعة بفرز \texttt{wakeup()} ولم تجد أي عملية نائمة بعد. بعد ذلك تنام العملية إلى الأبد.

لتجنب ذلك، تتطلب \texttt{sleep()} حيازة قفل الشرط (Condition Lock) الممرر إليها:
تقوم \texttt{sleep()} بحيازة قفل العملية الداخلي \texttt{p->lock} ثم تحرير قفل الشرط الممرر بشكل ذري قبل دعوة \texttt{sched()}.

\section{الكود البرمجي: النوم واليقظة}
\label{sec:code_sleep_wakeup}

تنفذ الدالة \texttt{sleep()} في \texttt{kernel/proc.c} بالخطوات المحددة التالية:
\begin{lstlisting}
void
sleep(void *chan, struct spinlock *lk)
{
  struct proc *p = myproc();
  
  // حيازة p->lock لحماية حالة العملية
  acquire(&p->lock);
  
  // تحرير قفل الشرط الممرر
  release(lk);

  // وضع القناة وتعديل الحالة إلى SLEEPING
  p->chan = chan;
  p->state = SLEEPING;

  sched();

  // بعد الاستيقاظ: تنظيف القناة وتحرير p->lock
  p->chan = 0;
  release(&p->lock);

  // إعادة حيازة قفل الشرط الأصلي
  acquire(lk);
}
\end{lstlisting}

وتقوم الدالة \texttt{wakeup(chan)} بتمشيط جدول العمليات بحثًا عن العمليات الممتلكة لـ \texttt{p->state == SLEEPING} والمطابقة لـ \texttt{p->chan == chan}، وتغير حالتها فورًا إلى \texttt{RUNNABLE}.

\section{الكود البرمجي: الأنابيب}
\label{sec:code_pipes_sleep}

تستخدم الأنابيب آلية \texttt{sleep()} و \texttt{wakeup()} للتنسيق بين الكاتب والقارئ:
إذا امتلاء الأنبوب، ينام الكاتب على قناة الكتابة (\texttt{\&p->nwrite}). وعندما يقرأ القارئ بيانات، يوقظ الكاتب عبر \texttt{wakeup(\&p->nwrite)}.

\section{الكود البرمجي: wait و exit و kill}
\label{sec:code_wait_exit_kill}

تتفاعل آليات النوم واليقظة لدعم دورة حياة العمليات:
\begin{itemize}
  \item \texttt{exit(status)}: تُستدعى لتدمير العملية الحالية. تغير حالة العملية إلى \texttt{ZOMBIE}، وتمرر أطفال العملية إلى العملية الأولى (\texttt{init})، وتوقظ الأب المنتظر عبر \texttt{wakeup(p->parent)} ثم تدعو \texttt{sched()} دون العودة مطلقًا.
  \item \texttt{wait(status)}: تبحث عن أي عملية ابن حالتها \texttt{ZOMBIE}. إذا وجدتها، تنظف ذاكرة الابن ومكدسه وإطار مصيدته وترجع معرّف الابن \texttt{pid}. وإذا كان للعملية أبناء ما زالوا يشتغلون، تنام العملية على عنوانها الذاتي \texttt{sleep(p, \&p->lock)} حتى ينهي أحدهم.
  \item \texttt{kill(pid)}: تعين الراية \texttt{p->killed = 1} للعملية المستهدفة وتوقظها فورًا إذا كانت نائمة. عند خروج العملية من المصيدة أو استدعاء النظام التالي، تفحص الراية وتنهي تنفيذها فورًا عبر \texttt{exit(-1)}.
\end{itemize}

\section{حماية وتأمين العمليات}
\label{sec:process_locking}

تستخدم النواة القفل المخصص لكل عملية (\texttt{p->lock}) لحماية التغييرات في \texttt{p->state}، \texttt{p->chan}، و \texttt{p->killed} لمنع تداخل الاستدعاءات المتوازية بين المعالجات المختلفة.

\section{العالم الحقيقي}
\label{sec:real_world_sleep}

تُعرف آليات النوم واليقظة في أدبيات أنظمة التشغيل بـ \textbf{المراقِبات} (Monitors) أو متغبرات الشرط (Condition Variables).
تستخدم الأنظمة الإنتاجية هياكل بيانات متقدمة تُدعى سلاسل الانتظار (Wait Queues) لربط القنوات بقوائم انتظار سريعة لتجنب المرور على جميع العمليات في النظام كما يكتفي \texttt{xv6}.

\section{تمارين}
\label{sec:exercises_sleep}

\begin{enumerate}
  \item اشرح بالتفصيل كيف تمنع حيازة \texttt{p->lock} وقفل الشرط حدوث مشكلة «الاستيقاظ المفقود».
  \item قم بتعديل \texttt{wakeup()} في \texttt{xv6} لتستند على مصفوفة قوائم انتظار لكل قناة نوم لزيادة السرعة.
\end{enumerate}
